%%%%%%%%%%%%%%%%%%%%%%%%%%%%%%%%%%%%%%%%%%%%%%%%%%%%%%%%%%%%%%%%%%%%%%%%%%%%%%%%
% LateX
%
% Author: FOSSEE 
% Copyright (c) 2009, FOSSEE, IIT Bombay
%%%%%%%%%%%%%%%%%%%%%%%%%%%%%%%%%%%%%%%%%%%%%%%%%%%%%%%%%%%%%%%%%%%%%%%%%%%%%%%%

\documentclass[14pt,compress]{beamer}

\mode<presentation>
{
  \usetheme{Warsaw}
  \useoutertheme{infolines}
  \setbeamercovered{transparent}
}

\usepackage[english]{babel}
\usepackage[latin1]{inputenc}
%\usepackage{times}
\usepackage[T1]{fontenc}

% Taken from Fernando's slides.
\usepackage{ae,aecompl}
\usepackage{mathpazo,courier,euler}
\usepackage[scaled=.95]{helvet}

\definecolor{darkgreen}{rgb}{0,0.5,0}

\usepackage{listings}
\lstset{language=bash,
    basicstyle=\ttfamily\bfseries,
    commentstyle=\color{red}\itshape,
  stringstyle=\color{darkgreen},
  showstringspaces=false,
  keywordstyle=\color{blue}\bfseries}

\newcommand{\inctime}[1]{\addtocounter{time}{#1}{\tiny \thetime\ m}}

\newcommand{\typ}[1]{\lstinline{#1}}

\newcommand{\kwrd}[1]{ \texttt{\textbf{\color{blue}{#1}}}  }

\setbeamercolor{emphbar}{bg=blue!20, fg=black}
\newcommand{\emphbar}[1]


\begin{document}

\begin{frame}

\begin{center}
\vspace{12pt}
\textcolor{blue}{\huge {\LaTeX}: Typesetting Text}
\end{center}
\vspace{18pt}
\begin{center}
\vspace{10pt}
\includegraphics[scale=0.95]{../images/fossee-logo.png}\\
\vspace{5pt}
\scriptsize Developed by FOSSEE Team, IIT-Bombay. \\ 
\scriptsize Funded by National Mission on Education through ICT\\
\scriptsize  MHRD,Govt. of India\\
\includegraphics[scale=0.15]{../images/iitb-logo.jpg}\\
\end{center}
\end{frame}

\begin{frame}
  \frametitle{Objectives}
  At the end of this session, you will be able to:
  \begin{itemize}
  \item Learn how to typeset your document using {\LaTeX}
  \item Use lists, listings in your document for formatting text
  \end{itemize}
\end{frame}

\begin{frame}
  \frametitle{Prerequisite}
  Spoken Tutorial on:
  \begin{itemize}
  \item {\LaTeX} Part 1 - Getting Started
  \item {\LaTeX} Part 2 - Introduction
  \item {\LaTeX} Part 3 - Structuring the Content
  \end{itemize}
\end{frame}

\begin{frame}[fragile]
  \frametitle{Quotation Marks}
  \begin{itemize}
  \item Use \`~ (accent) for left quote
  \item Use \'~ (apostrophe) for right quote
  \item For double quotes, use them twice
  \end{itemize}
  \tiny See rev11 of \typ{hg}
\end{frame}

\begin{frame}[fragile]
  \frametitle{Fonts - Emphasis, Fixed width, \ldots}
  \begin{itemize}
  \item \lstinline{\emph} gives emphasized or italic text
  \item \typ{flushleft} to have text left aligned
  \item \typ{flushright}, \typ{center}
  \end{itemize}
  \tiny See rev12 of \typ{hg}
\end{frame}

\begin{frame}[fragile]
  \frametitle{Fonts - Emphasis, Fixed width, \ldots}
  \begin{itemize}
  \item \lstinline{\texttt} gives fixed width font
  \item \lstinline{\textbf} bold face font
  \item \lstinline{--} en dash (--); \lstinline{---} em dash (---). 
  \end{itemize}
  \tiny See rev13 of \typ{hg}
\end{frame}

\begin{frame}[fragile]
  \frametitle{Lists}
  \begin{itemize}
  \item \lstinline{enumerate} environment is used for numbered lists
  \item \lstinline{itemize} environment gives un-numbered lists
  \item Each item in the list is specified using \lstinline{\item}
  \item Nested lists are also easily handled, as expected
  \end{itemize}
  \tiny See rev14 of \typ{hg}
\end{frame}

\begin{frame}[fragile]
  \frametitle{Footnotes}
  \begin{itemize}
  \item \typ{\\footnote} command adds a footnote
  \end{itemize}
  \tiny See rev15 of \typ{hg}
\end{frame}

\begin{frame}[fragile]
  \frametitle{Labels and References}
  \begin{itemize}
  \item \lstinline+\label{labelname}+ is used to label an element
  \item \lstinline+\ref{labelname}+ is used to refer to that element
  \item Compile twice
  \end{itemize}
  \tiny See rev15 of \typ{hg}
\end{frame}

\begin{frame}[fragile]
  \frametitle{Including code}
  \begin{itemize}
  \item Instead of using \lstinline{\texttt} we could use
    \lstinline{\verbatim} 
  \item \lstinline+listings+ is a powerful package
  \item \lstinline+\usepackage{listings}+ needs to be added 
  \item Tell {\LaTeX} the language to be used, using \typ{\\lstset}
  \end{itemize}
  \tiny See rev16 of \typ{hg}
\end{frame}

\begin{frame}[fragile]
  \frametitle{Including code}
  \begin{itemize}
  \item Use \lstinline+\lstlisting+ for a block of code
  \item \typ{\\lstinline} for inline code
  \end{itemize}
  \tiny See rev16 of \typ{hg}
\end{frame}


\begin{frame}[fragile]
	\frametitle{Summary...}
	\begin{itemize}
	\item Put Quotation Marks around text
	\item How to Emphasize and give fixed width to fonts.
	\item Use numbered and un-numbered lists
	\item Add Footnotes, Labels and References
	\item Use the listings package to include code
	\end{itemize}
\end{frame}

\begin{frame}[fragile]
\frametitle{Evaluation}
\begin{enumerate}
\item Which environment is used including a block of code?
\item What are the basic lists that make environments?
\item Joe has numerous used labels inside his Latex document. But all the references to label names come up as question marks. What might be the problem? 
\end{enumerate}
\end{frame}
\begin{frame}

\frametitle{Solutions}
\begin{enumerate}
\item lstlistings
\vspace{15pt}
\item
\begin{enumerate} 
\item itemize(for unorderd lists) 
\item enumerate(for ordered lists)
\item description(lists where you want to specify your own label)
\end{enumerate}
\item While using labels, the latex document should be compiled twice for the references to show up.
\end{enumerate}
\end{frame}
\begin{frame}

\begin{block}{}
  \begin{center}
  \textcolor{blue}{\Large THANK YOU!} 
  \end{center}
  \end{block}
\begin{block}{}
  \begin{center}
    For more Information, visit our website\\
    \url{http://fossee.in/}
  \end{center}  
  \end{block}
\end{frame}

\end{document}
